\subsubsection*{An integrality lemma for cyclotomic algebraic bases}

Let $k\subset\overline{\mathbb Q}$ be contained in
$F\mathbb Q^{\mathrm{ab}}$ for a number field $F$, and let $L/k$ be
finite. Let $b_1,\ldots,b_r\in L$ be distinct and put
$V=\mathbb A^1_L\setminus\{b_1,\ldots,b_r\}$.
Suppose
\[
 \xi\in\varprojlim_n H^1_{\mathrm{et}}(V,\mu_n)
\]
has the following evaluation property: for every $a\in V(L)$,
$a^*\xi$ is the Kummer class of an element of $L^*$.
Then $\xi$ is the Kummer class of an actual unit of $V$.
This is an integrality assertion for a cohomology class; the individual
elements realizing its evaluations need not initially form a rational
function.

\begin{proof}
Since $\operatorname{Pic}(V)=0$ and
$\mathcal O(V)^*=L^*\times\prod_i(t-b_i)^{\mathbb Z}$, Kummer theory gives
\[
 \xi=\gamma+\sum_{i=1}^r\lambda_i[t-b_i],\qquad
 \gamma\in\widehat{L^*},\quad \lambda_i\in\widehat{\mathbb Z}.
                                                        \tag{CI.1}
\]
All completions in this argument are $\varprojlim_n A/nA$ (in
multiplicative notation, $\varprojlim_n A/A^n$). We first prove
$\lambda_i\in\mathbb Z$ without assuming that $L$ has a discrete
valuation.

Choose a number field $K_0\subset L$ containing all the $b_i$.
For each $i$, choose a different prime ideal $\mathfrak p_i$ of
$\mathcal O_{K_0}$ such that the $b_j$ are integral there, their pairwise
differences are units there, and the residue field has more than $r$
elements. Only finitely many primes are excluded. The Chinese remainder
theorem supplies $x,y\in K_0$ such that
\[
 v_{\mathfrak p_i}(x-b_i)=1,\qquad
 v_{\mathfrak p_i}(x-b_j)=0\ (j\ne i),\qquad
 v_{\mathfrak p_i}(y-b_j)=0\ (\text{all }j).
                                                        \tag{CI.2}
\]
For the first condition specify $x$ modulo $\mathfrak p_i^2$; for the
last choose a residue different from each $b_j$. Thus $x,y\in V(L)$.
The elements $r_j=(x-b_j)/(y-b_j)$ are multiplicatively independent
modulo roots of unity: applying the valuations in (CI.2) to a relation
$\prod_j r_j^{e_j}\in\mu_\infty$ gives $e_i=0$ for every $i$.
This independence persists in every algebraic overfield, since the
relation itself is an equality of the same algebraic numbers.

There is a number field $F_0$ for which $L\subset M=F_0\mathbb
Q^{\mathrm{ab}}$: take $F_0$ to contain $F$ and finitely many algebraic
generators of $L/k$. Theorem~1.1 of Amoroso and Zannier (2000) implies
\[
 \bigcap_{n\ge1}(M^*)^n=\mu_\infty.
                                                        \tag{CI.3}
\]
Indeed $[\mathbb Q^{\mathrm{ab}}(a):\mathbb Q^{\mathrm{ab}}]
\le [M:\mathbb Q^{\mathrm{ab}}]$ for $a\in M$, so their height bound
is uniformly positive on the non-torsion elements of $M$. If
$z_n^n=z\ne0$ for every $n$, then $h(z_n)=h(z)/n$, forcing $z$ to be
a root of unity. Conversely every root of unity is divisible in $M$,
which contains all roots of unity. Tsujimura (2022,
Proposition~4.9), applied with its parameter $f=1$ and number field
$F_0$, now says that
\[
 A=M^*/\mu_\infty
\]
is a free abelian group. This is the exact specialization of that
theorem used here.

The independent elements $[r_1],\ldots,[r_r]$ of $A$ have finite
support in a basis of the free abelian group. Their coordinate matrix
has rank $r$, so we may choose homomorphisms
$\chi_1,\ldots,\chi_r:A\to\mathbb Z$ such that the integer matrix
$N=(\chi_i([r_j]))_{i,j}$ has nonzero determinant. Choose actual
elements $u_x,u_y\in L^*$ realizing the two evaluations of $\xi$.
Subtracting the two evaluations of (CI.1), then mapping to
$\widehat A$ and applying the completed characters, gives
\[
 N(\lambda_1,\ldots,\lambda_r)^{\mathrm t}
   =\bigl(\chi_i([u_x/u_y])\bigr)_i\in\mathbb Z^r.
                                                        \tag{CI.4}
\]
The adjugate matrix shows that
$d\lambda_i\in\mathbb Z$, where $d=\det N\ne0$.
For each prime dividing $d$, membership in $\mathbb Z_p$ forces the
corresponding divisibility of that integer by the prime power in $d$.
Consequently $\lambda_i\in\mathbb Z$. Equivalently,
$\mathbb Q\cap\widehat{\mathbb Z}=\mathbb Z$ inside the finite
adeles. Put $m_i=\lambda_i$.

Finally choose any $a\in V(L)$ and an element $u_a\in L^*$ realizing
$a^*\xi$. The constant class in (CI.1) is then the Kummer class of
$c=u_a/\prod_i(a-b_i)^{m_i}$. Thus
$\xi=[c\prod_i(t-b_i)^{m_i}]$, as required. When $r=0$ this last
evaluation argument proves the result directly. The actual unit is
unique modulo $\bigcap_n(L^*)^n$, a subgroup of the roots of unity;
the assertion does not identify different pointwise choices before
the translated-function argument is applied.
\end{proof}
